% Generated from locale/tr/content/incompleteness/representability-in-q/c.tex; do not hand-edit.


\olfileid[tr]{inc}{req}{cee}
\olsection{$C$ Fonksiyonları}

$C$, aşağıdakileri içeren en küçük fonksiyonlar kümesi olsun:
\begin{enumerate}
\item $0$,
\item ardıl,
\item toplama,
\item çarpma,
\item izdüşümler ve
\item eşitliğin karakteristik fonksiyonu $\Char{=}$;
\end{enumerate}
ayrıca şu işlemler altında kapalı olsun:
\begin{enumerate}
\item bileşke ve
\item düzenli fonksiyonlara uygulanan sınırsız arama.
\end{enumerate}
Son kısıtın yalnızca, sonucu tümel olduğunda $\mu$ işlemini
kullanabileceğiniz anlamına geldiğini hatırlayın. Bunu \emph{genel
özyinelemeli} fonksiyonların tanımıyla karşılaştırın: burada toplama,
çarpma ve $\Char{=}$ eklenmiş, fakat ilkel özyineleme çıkarılmıştır.

Toplama, çarpma ve $\Char{=}$ özyinelemeli olduğundan $C$ içindeki her şey
açıkça özyinelemelidir. Tersinin de doğru olduğunu göstereceğiz; bu, $C$
içindeki öteki araçlarla ilkel özyinelemeyi gerçekleştirebileceğimizi
söylemeye varır.

